\documentclass[10pt,a4paper]{amsart}
\usepackage[margin=19mm]{geometry}
\usepackage{amsmath,amssymb,mathtools}
\usepackage[scr=rsfs,cal=euler]{mathalfa}
\usepackage{xfrac}
\usepackage[unicode,hidelinks,bookmarks=false]{hyperref}
\newcommand{\ie}{i.e.\ }
\newcommand{\cf}{cf.\ }
\newcommand{\resp}{respectively\ }
\newcommand{\Subsection}{\subsection}
\providecommand{\nobreakdash}{\nobreak\hbox{-}}
\setlength{\emergencystretch}{2em}
\multlinegap0pt
\usepackage[all]{xy}
\newtheorem{thm}{Theorem}
\newtheorem{lem}{Lemma}
\newtheorem{dfn}{Definition}
\newtheorem*{nota}{Notations}
\newtheorem*{claim}{Claim}
\newcommand{\Z}{\mathbb{Z}}
\newcommand{\Zpn}{\mathbb{Z}_{>0}}
\newcommand{\cH}{\mathcal{H}}
\DeclareMathOperator{\Ker}{Ker}
\DeclareMathOperator{\Ima}{Im}
\DeclareMathOperator{\Coker}{Coker}
\DeclareMathOperator{\Hom}{Hom}
\DeclareMathOperator{\Ind}{Ind}
\DeclareMathOperator{\pr}{pr}
\DeclareMathOperator{\rk}{rank}
\DeclareMathOperator{\tor}{tor}
\title[The quasi-permutation theorem for $G_{m,\nu}=C_{m}\times D_{2^{\nu}}$]{The rationality problem for multinorm one tori: the quasi-permutation theorem for $G_{m,\nu}=C_{m}\times D_{2^{\nu}}$ (Theorem 6.7)}
\author{Sumito Hasegawa, Kazuki Kanai, Yasuhiro Oki}
\date{Source: arXiv:2504.04078v4 (CC BY 4.0)}
\begin{document}
\begin{abstract}
In this paper, we study the rationality problem for multinorm one tori, a natural generalization of norm one tori. 
For multinorm one tori that split over finite Galois extensions with nilpotent Galois group, we prove that stable rationality and retract rationality are equivalent, and give a criterion for the validity of the above two conditions. 
This generalizes the result of Endo (2011) on the rationality problem for norm one tori. 
To accomplish it, 
we introduce a generalization of character groups of multinorm one tori.
Moreover, we establish systematic reduction methods originating in work of Endo (2001) for an investigation of the rationality problem for arbitrary multinorm one tori.
In addition, we provide a new example for which the multinorm principle holds.


\end{abstract}
\maketitle

\noindent\emph{Source and licence.} The rationality problem for multinorm one tori. Version arXiv:2504.04078v4 (CC BY 4.0), distributed under the Creative Commons Attribution 4.0 International licence, as recorded in the arXiv licence metadata for this exact version and shown on the abstract page https://arxiv.org/abs/2504.04078v4. Full rights notice: licenses/arxiv-2504.04078-CC-BY-4.0.txt.\par\smallskip

\noindent\textit{Editorial scope.} Counted unit: the theorem printed as Theorem 6.7 of the article (source0.tex lines 1942--1949, source label \texttt{thm:nlqp}), delivered together with the article's own complete original proof of it (source0.tex lines 1951--2132, one \texttt{proof} environment of 182 lines that proves the whole assertion with no gap and no deferral). No mathematics is written, repaired or re-derived: statement and proof are the article's own text line for line, in the article's own notation ($G_{m,\nu}$, $\cH_{m,\nu}$, $J_{G/\cH}$, $I_{G/\cH}$, $U_{m,\nu}$, $R_{m,\nu}$, $\Phi_{m,\nu}$, $\Xi_{m,\nu}$, $S_{m,\nu}$, $Z_{\nu}$). Required context retained uncounted: the article's abstract (lines 197--203); its notations block defining the normal core and the normalizer, $G$-lattices, the dual lattice $M^{\circ}$ and the lattices $M^{N}$ and $M^{[N]}$ (lines 490--515); its definition of permutation, quasi-permutation and quasi-invertible lattices (lines 535--546, cited to Endo); its definition of stably permutation, invertible, coflabby and flabby lattices (lines 554--562, cited to Lorenz); its definition of coflabby and flabby resolutions (lines 581--595, cited to Lorenz); its construction of the augmentation maps, of the induced maps and of the lattices $I_{G/\cH}$ and $J_{G/\cH}$ (lines 786--810); its notation for the dihedral group $D_{n}$ (lines 1717--1720); the statement of the article's Theorem 6.1 (lines 1722--1732), which the counted proof invokes twice; the statements of the article's Lemmas 6.3, 6.4, 6.5 and 6.6 (lines 1759--1778, 1795--1810, 1820--1834 and 1848--1888), each of which the counted proof invokes by label; and the article's own sentence announcing the counted theorem (line 1940). Every definition, numbered display and label of those lines is retained unchanged. Omitted from this package and not redistributed: the article's preamble and title page (lines 1--196); the remainder of its front matter with the table of contents (lines 204--209); its whole introduction, containing the announcements of its five main theorems and its acknowledgments (lines 210--489); everything lying between the retained excerpts of its second and third sections (lines 516--534, 547--553, 563--580, 596--785); the rest of its third section with the weighted lattices $I^{(\varphi)}$ and $J^{(\varphi)}$ and its two reduction subsections (lines 811--1716); its fourth and fifth sections on $p$-groups and on the proof of its Theorem 1.1 (lines 1470--1712); the parts of its sixth section that lie between the retained excerpts and its second subsection on non-quasi-invertible lattices (lines 1721, 1733--1758, 1779--1794, 1811--1819, 1835--1847, 1889--1939, 1941, 1950, 2133--2280); its seventh, eighth and ninth sections containing the proofs of its Theorems 1.2, 1.3 and 1.5 (lines 2281--2927); its bibliography (lines 2928--3013); and its closing end-document line (line 3015). Nothing inside a retained excerpt is omitted or altered: every retained body is delivered exactly as the article's lines are written, so no omission receipt of any kind accompanies this package. Citations: the retained excerpts carry exactly three citation commands, all of them in the optional arguments of the three retained definition headers, and together they cite exactly two keys of the article's own bibliography, \texttt{Endo2011} and \texttt{Lorenz2005}. Those two entries are transcribed verbatim from the article's own in-source \texttt{thebibliography} (lines 2928--3013, 41 entries) in the article's own order and with the article's own printed labels, so the delivered citation marks read [End11] and [Lor05] as the article's do. No other entry of the article's bibliography is transcribed or redistributed. Reference adaptation: none was needed and none was made. Every reference inside the delivered unit resolves inside it: the labels \texttt{thm:dnqp}, \texttt{lem:cfad}, \texttt{lem:lmis}, \texttt{lem:fgdc} and \texttt{lem:cssp} of the retained statements, which the counted proof cites by name; the internal labels \texttt{eq:gmnr}, \texttt{eq:gmna}, \texttt{eq:fcua}, \texttt{eq:fmns}, \texttt{eq:xitc}, \texttt{eq:fsdd}, \texttt{eq:imxi}, \texttt{eq:ffex} and \texttt{eq:aflm} of the numbered displays inside the counted proof, together with \texttt{eq:adex} inside the retained Lemma 6.3; and the label \texttt{thm:nlqp} of the counted statement itself. No unresolved reference or citation is produced. Printed slips kept literal: no printed word, symbol, hypothesis or number of the counted statement or of its proof was altered. In particular the article's announcing sentence calls the counted theorem's parameter $n$ although the theorem itself uses $\nu$, exactly as the article prints it. Layout and class: the article's own class is \texttt{amsart} with the option 12pt, and the article uses no publisher class or style file, so no class or style file travels with this package and none is redistributed. The wrapper is the lane's pinned \texttt{amsart} editorial wrapper at 10pt on A4, which adds the article's own macro definitions that the retained text needs (the blackboard $\Z$, the positive integers $\Zpn$, the calligraphic $\cH$, and the operators $\Ker$, $\Ima$, $\Coker$, $\Hom$, $\Ind$, $\pr$, $\rk$ and $\tor$), declares the theorem-like environments the retained text needs (\texttt{thm}, \texttt{lem}, \texttt{dfn}, \texttt{nota} and \texttt{claim}), and loads the \texttt{xy} package that the retained commutative diagrams need. The delivered numbering is the unit's own: the retained definitions print as Definitions 1--3, the retained Theorem 6.1 as Theorem 1, the retained Lemmas 6.3--6.6 as Lemmas 1--4, the counted Theorem 6.7 as Theorem 2, and the retained Notations and Claim print unnumbered. The article's 12pt measure differs from the wrapper's 10pt; that is a page-appearance difference of the wrapper only, and no formula, symbol, hypothesis or argument changes. Assets: the article has no figure environment, no graphics-inclusion command, no PDF-page inclusion, no table and no TikZ picture anywhere in its source, and no retained span contains a graphics-inclusion command, an input directive or an include directive; the package declares no asset dependency and no image file is copied into or redistributed with it, the only diagram code retained being the article's own \texttt{xymatrix} source, which the \texttt{xy} package renders from text. Licence: version arXiv:2504.04078v4 of this article is distributed under the Creative Commons Attribution 4.0 International licence, as recorded in the arXiv licence metadata for this exact version and shown on the abstract page https://arxiv.org/abs/2504.04078v4. A full rights notice is delivered at licenses/arxiv-2504.04078-CC-BY-4.0.txt. Characters: every retained excerpt is pure ASCII, so the delivered engine, XeLaTeX with its default Latin Modern text font, reports no missing character. No glyph is substituted and no word is rewritten.
\begin{nota}
Let $G$ be a finite group. 
\begin{itemize}
\item For a subgroup $H$ of $G$, we define $N^{G}(H)$ and $N_{G}(H)$ as follows: 
\begin{equation*}
N^{G}(H):=\bigcap_{g\in G}gHg^{-1},\quad
N_{G}(H):=\{g\in G\mid gHg^{-1}=H\}. 
\end{equation*}
Note that $N^{G}(H)$ and $N_{G}(H)$ are called the \emph{normal core} of $H$ and the \emph{normalizer} of $H$ in $G$, respectively. 
\item A \emph{$G$-lattice} refers to a finitely generated free abelian group equipped with a left action of $G$. For a $G$-lattice $M$, the dual lattice of $M$ is denoted by
\begin{equation*}
    M^{\circ}:=\Hom_{\Z}(M,\Z). 
\end{equation*}
Here we define a left action of $G$ on $M^{\circ}$ as
\begin{equation*}
    G\times M^{\circ} \rightarrow M^{\circ};\,(g,f)\mapsto [x\mapsto f(g^{-1}x)]. 
\end{equation*}
%\item For a $G$-module $M$, we write $M_{\tor}$ for the torsion part of $M$. 
\item Let $M$ be a $G$-lattice. For a normal subgroup $N$ of $G$, we define $G/N$-lattices $M^{N}$ and $M^{[N]}$ as follows: 
\begin{equation*}
M^{N}:=\{x\in M\mid n(x)=x\text{ for all }n\in N\},\quad 
M^{[N]}:=\left((M^{\circ})^{N}\right)^{\circ}. 
\end{equation*}
Note that $M^{[N]}$ is isomorphic to $M_{N}/M_{N,\tor}$, where $M_{N}$ is the $N$-coinvariant part of $M$, and $M_{N,\tor}$ is the torsion part of $M_{N}$. Note that $M^{N}$ and $M^{[N]}$ may not coincide in general. 
\end{itemize}
\end{nota}

\begin{dfn}[{\cite[\S 1]{Endo2011}}]
Let $G$ be a finite group. We say that a $G$-lattice $M$ is
\begin{enumerate}
    \item \emph{permutation} if $M$ has a $\Z$-basis permuted by $G$, that is, $M\cong \bigoplus_{i=1}^{m}\Z[G/H_i]$ for some subgroups $H_1,H_2,\dots, H_m$;
    \item \emph{quasi-permutation} if there is an exact sequence of $G$-lattices
    \begin{equation*}
        0\rightarrow M \rightarrow R \rightarrow F \rightarrow 0, 
    \end{equation*}
    where $R$ and $F$ are permutation;
    \item \emph{quasi-invertible} if it is a direct summand of a quasi-permutation $G$-lattice. 
\end{enumerate}
\end{dfn}

\begin{dfn}[{\cite[\S\S 2.3--2.5]{Lorenz2005}}]
Let $G$ be a finite group. We say that a $G$-lattice $M$ is
\begin{enumerate}
\item \emph{stably permutation} if $M\oplus R\cong R'$ for some permutation $G$-lattices $R$ and $R'$;
\item \emph{invertible} (or, \emph{permutation projective}) if it is a direct summand of a permutation $G$-lattice;
\item \emph{coflabby} if $H^{1}(H,M)=0$ for any subgroup $H$ of $G$;
\item \emph{flabby} if $M^{\circ}$ is coflabby. 
\end{enumerate}
\end{dfn}

\begin{dfn}[{\cite[\S 2.6]{Lorenz2005}}]
Let $G$ be a finite group, and $M$ a $G$-lattice. 
\begin{enumerate}
\item A \emph{coflabby resolution} of $M$ is an exact sequence of $G$-lattices
\begin{equation*}
0\rightarrow U \rightarrow R \rightarrow M \rightarrow 0, 
\end{equation*}
where $R$ is permutation and $U$ is coflabby. 
\item A \emph{flabby resolution} of $M$ is an exact sequence of $G$-lattices
\begin{equation*}
0\rightarrow M \rightarrow R \rightarrow F \rightarrow 0, 
\end{equation*}
where $R$ is permutation and $F$ is flabby. 
\end{enumerate}
\end{dfn}

Let $G$ be a finite group, and $H$ a subgroup of $G$. Then one has a surjection
\begin{equation*}
    \varepsilon_{G/H}\colon \Z[G/H]\rightarrow \Z;\,\sum_{g \in G/H} a_g g\mapsto \sum_{g \in G/H} a_g, 
\end{equation*}
which is called the \emph{augmentation map}. Moreover, the dual of $\varepsilon_{G/H}$ coincides with the homomorphism
\begin{equation*}
    \varepsilon_{G/H}^{\circ}\colon \Z\rightarrow \Z[G/H];\,1\mapsto \sum_{g\in G/H}g. 
\end{equation*}
On the other hand, consider a finite group $G$ and subgroups $H'\subset H$ of $G$. Then the homomorphisms of $H$-lattices
\begin{equation*}
   \varepsilon_{H/H'}\colon \Z[H/H']\rightarrow \Z,\quad \varepsilon_{H/H'}^{\circ}\colon \Z \rightarrow \Z[H/H']
\end{equation*}
induce homomorphisms of $G$-lattices
\begin{equation*}
   \Ind_{H}^{G}\varepsilon_{H/H'}\colon \Z[G/H']\rightarrow \Z[G/H],\quad \Ind_{H}^{G}\varepsilon_{H/H'}^{\circ}\colon \Z[G/H] \rightarrow \Z[G/H']. 
\end{equation*}

For a multiset $\cH$ of subgroups of $G$, we define a $G$-module $I_{G/\cH}$ by an exact sequence
\begin{equation*}
    0\rightarrow I_{G/\cH}\rightarrow \bigoplus_{H\in \cH}\Z[G/H]\xrightarrow{(\varepsilon_{G/H})_{H\in \cH}} \Z \rightarrow 0. 
\end{equation*}
Furthermore, we define $J_{G/\cH}:=I_{G/\cH}^{\circ}$. Then one has an exact sequence
\begin{equation*}
    0\rightarrow \Z \xrightarrow{(\varepsilon_{G/H}^{\circ})_{H\in \cH}} \bigoplus_{H\in \cH}\Z[G/H] \rightarrow J_{G/\cH} \rightarrow 0. 
\end{equation*}

Here we give some examples of $G$-lattices $J_{G/\cH}^{(\varphi)}$ that are quasi-permutation. For $n\in \Zpn$, we denote by $D_{n}$ the dihedral group of order $2n$, that is, 
\begin{equation*}
D_{n}:=\langle \sigma_{n},\tau_{n} \mid \sigma_{n}^{n}=\tau_{n}^{2}=1,\tau_{n}\sigma_{n}\tau_{n}=\sigma_{n}^{-1}\rangle. 
\end{equation*}

\begin{thm}\label{thm:dnqp}
Let $m$ be a positive integer. Consider a strongly reduced set 
\begin{equation*}
\cH:=\{\langle \tau_{2m} \rangle, \langle \sigma_{2m}\tau_{2m} \rangle \}
\end{equation*}
of subgroups of $D_{2m}$. Then there is an exact sequence of $D_{2m}$-lattices 
\begin{equation*}
    0\rightarrow J_{D_{2m}/\cH}\rightarrow \Z[D_{2m}]\oplus \Z \rightarrow \Z[D_{2m}/\langle \sigma_{2m} \rangle]\rightarrow 0. 
\end{equation*}
In particular, $J_{D_{2m}/\cH}$ is quasi-permutation. 
\end{thm}

\begin{lem}\label{lem:cfad}
Let $G$ be a finite group, and 
\begin{equation*}
0\rightarrow F \xrightarrow{\iota} R \xrightarrow{\Phi} M \rightarrow 0
\end{equation*}
a coflabby resolution of a $G$-lattice $M$. Consider a homomorphism of $G$-lattices $\psi \colon R'\rightarrow M$, where $R'$ is a permutation $G$-lattice. We denote by $\Phi'$ the sum of $\Phi$ and $\psi$. Then there is an exact sequence
\begin{equation}\label{eq:adex}
0\rightarrow F\oplus R'\xrightarrow{\iota'} R\oplus R'\xrightarrow{\Phi'}M \rightarrow 0
\end{equation}
which satisfies a commutative diagram
\begin{equation*}
\xymatrix@C=35pt{
F\ar[r]^{x \mapsto (x,0)\hspace{10pt}}\ar[d]^{\iota}& F\oplus R'\ar[d]^{\iota'}&
R'\ar[l]_{\hspace{15pt}x\mapsto (0,x)}\ar@{=}[d]\\
R \ar[r]^{x\mapsto (x,0)\hspace{10pt}}& 
R\oplus R' \ar[r]^{\hspace{10pt}\pr_{2}}& R'. 
}
\end{equation*}
In particular, \eqref{eq:adex} is also a coflabby resolution of $M$. 
\end{lem}

\begin{lem}\label{lem:lmis}
Consider a commutative diagram of finite free abelian groups
\begin{equation*}
\xymatrix{
0\ar[r]& M_{1}\ar[r] \ar[d]^{f_{1}} & M_{2} \ar[r]\ar[d]^{f_{2}}& M_{3} \ar[r]\ar[d]^{f_{3}}& 0\\
0\ar[r]& M'_{1}\ar[r] & M'_{2} \ar[r] & M'_{3}, &
}
\end{equation*}
where the horizontal sequences are exact. We further assume that 
\begin{itemize}
    \item $\rk_{\Z}(M_{1})=\rk_{\Z}(M'_{1})$; 
    \item $f_{2}$ and $f_{3}$ are injective; and 
    \item the cokernel of $f_{2}$ is torsion-free. 
\end{itemize}
Then the homomorphism $f_{1}$ is an isomorphism. 
\end{lem}

\begin{lem}\label{lem:fgdc}
Let $G$ be a finite group, and $N$ its normal subgroup. Consider a commutative diagram of $G$-lattices as follows: 
\begin{equation*}
\xymatrix{
0\ar[r]& M_{1}\ar[r]\ar[d]^{f_{1}}& M_{2}\ar[r]\ar[d]^{f_{2}}& M_{3}\ar[r]\ar[d]^{f_{3}}&0\\
0\ar[r]& M'_{1}\ar[r]& M'_{2}\ar[r]& M'_{3}\ar[r]&0,  
}
\end{equation*}
where the horizontal sequences are exact. We further assume that 
\begin{itemize}
\item[(a)] $H^{1}(N,\Ker(f_{1}))=0$; and 
\item[(b)] $N$ acts trivially on $M_{3}$.  
\end{itemize}
Then $\Ker(f_{2})$ is generated by $\Ker(f_{1})$ and $\Ker(f_{2})^{N}$. 
\end{lem}

\begin{lem}\label{lem:cssp}
Let $G$ be a finite group, $H$ its subgroup, and $N_{1}$ and $N_{2}$ normal subgroups of $G$. Consider homomorphisms of $G$-lattices as follows: 
\begin{gather*}
f_{1}\colon M_{1}\oplus \Z[G/H]^{N_{1}}\rightarrow M'_{1}\oplus \Z[G/H];\\
f_{2}\colon M_{2}\oplus \Z[G/H]^{N_{2}}\rightarrow M'_{1}\oplus \Z[G/H]\oplus M'_{2}. 
\end{gather*}
Let $f$ be the sum of $f_{1}$ and $f_{2}$, and denote by $M' \subset M'_{1}\oplus \Z[G/H]\oplus M'_{2}$ the image of $f$. We further assume that
\begin{itemize}
\item[(1)] $\gcd((HN_{1}N_{2}:HN_{1}),(HN_{1}N_{2}:HN_{2}))=1$; 
\item[(2)] there exist commutative diagrams
\begin{equation*}
\xymatrix@C=35pt{
M_{1}\oplus \Z[G/H]^{N_{1}}\ar[d]^{f_{1}}&
\Z[G/H]^{N_{1}}\ar[l]_{\hspace{15pt}x\mapsto (0,x)}\ar[d]^{x\mapsto x}\\
M'_{1}\oplus \Z[G/H]\ar[r]^{\hspace{10pt}\pr_{2}}&\Z[G/H],
}\quad
\xymatrix@C=30pt{
M_{2}\oplus \Z[G/H]^{N_{2}}\ar[d]^{f_{2}}&
\Z[G/H]^{N_{2}}\ar[l]_{\hspace{15pt}x\mapsto (0,x)}\ar[d]^{x\mapsto x}\\
M'_{1}\oplus \Z[G/H]\oplus M'_{2}\ar[r]^{\hspace{25pt}\pr_{2}}&
\Z[G/H];
}
\end{equation*}
\item[(3)] $M'\cap (M'_{1}\oplus \{0\}\oplus M'_{2})^{N_{1}N_{2}}\subset f_{1}(M_{1}\oplus \{0\})+f_{2}(M_{2}\oplus \{0\})$; and 
\item[(4)] $\rk_{\Z}\Ker(f)=(G:HN_{1}N_{2})$. 
\end{itemize}
Then there is an isomorphism
\begin{equation*}
M' \oplus \Z[G/HN_{1}N_{2}]\cong M_{1}\oplus M_{2}\oplus \Z[G/HN_{1}] \oplus \Z[G/HN_{2}]
\end{equation*}
such that the following diagram is commutative: 
\begin{equation*}
\xymatrix{
M_{1}\oplus M_{2}\ar[r]^{f}\ar[d]^{(x,y)\mapsto (x,y,0,0)}&
M'\ar[d]^{x\mapsto (x,0)}\\
M_{1}\oplus M_{2}\oplus \Z[G/HN_{1}] \oplus \Z[G/HN_{2}]\ar[r]^{\hspace{40pt}\cong}&
M'\oplus \Z[G/HN_{1}N_{2}]. 
}
\end{equation*}
%$M_{1}\oplus M_{2}\oplus \{0\}\oplus \{0\}$ corresponds to a $G$-sublattice of $M'\oplus \{0\}$. 
\end{lem}

The following generalizes Theorem \ref{thm:dnqp} in the case that $n$ is a power of $2$.

\begin{thm}\label{thm:nlqp}
Let $m$ be an odd positive integer, and $\nu \in \Zpn$. Consider the finite group
\begin{align*}
G_{m,\nu}:&=C_{m}\times D_{2^{\nu}}\\
&=\langle \rho_{m},\sigma_{2^{\nu}},\tau_{2^{\nu}} \mid \rho_{m}^{m}=\sigma_{2^{\nu}}^{2^{\nu}}=\tau_{2^{\nu}}^{2}=1,\rho_{m}\sigma_{2^{\nu}}=\sigma_{2^{\nu}}\rho_{m},\rho_{m}\tau_{2^{\nu}}=\tau_{2^{\nu}}\rho_{m},\tau_{2^{\nu}}\sigma_{2^{\nu}}\tau_{2^{\nu}}^{-1}=\sigma_{2^{\nu}}^{-1} \rangle, 
\end{align*}
and define $\cH_{m,\nu}:=\{\langle \tau_{2^{\nu}} \rangle,\langle \sigma_{2^{\nu}}\tau_{2^{\nu}} \rangle \}$. Then the $G_{m,\nu}$-lattice $J_{G_{m,\nu}/\cH_{m,\nu}}$ is quasi-permutation. 
\end{thm}

\begin{proof}
In this proof, put $I_{m,\nu}:=I_{G_{m,\nu}/\cH_{m,\nu}}$. By Theorem \ref{thm:dnqp}, we may assume $m>1$. Consider homomorphisms of $G$-lattices as follows: 
\begin{gather*}
\psi_{m,\nu,0}\colon R_{m,\nu,0}:=\Z[G_{m,\nu}/\langle \rho_{m} \rangle]\rightarrow I_{m,\nu};\,1\mapsto \left(\sum_{i=0}^{m-1}\rho_{m}^{i},-\sum_{i=0}^{m-1}\rho_{m}^{i} \right);\\
\psi_{m,\nu,1}\colon R_{m,\nu,1}:=\Z \rightarrow I_{m,\nu};\,1\mapsto \left(\left(\sum_{i=0}^{m-1}\rho_{m}^{i}\right)\left(\sum_{j=0}^{2^{\nu}-1}\sigma_{2^{\nu}}^{j}\right),-\left(\sum_{i=0}^{m-1}\rho_{m}^{i}\right)\left(\sum_{j=0}^{2^{\nu}-1}\sigma_{2^{\nu}}^{j}\right) \right);\\
\psi_{m,\nu,2}\colon R_{m,\nu,2}:=\Z[G_{m,\nu}/\langle \tau_{2^{\nu}} \rangle]\rightarrow I_{m,\nu};\,1\mapsto (1+\rho_{m},-(1+\sigma_{2^{\nu}}^{-1}));\\
\omega_{m,\nu,i}\colon R'_{m,\nu,i}:=\Z[G_{m,\nu}/\langle \sigma_{2^{\nu}}^{2^{\nu-i}}, \sigma_{2^{\nu}} \tau_{2^{\nu}} \rangle]\rightarrow I_{m,\nu};\,1\mapsto \left(0,\left(\sum_{j=0}^{2^{i}-1}\sigma_{2^{\nu}}^{2^{\nu-i}j}\right)(1-\rho_{m}\sigma_{2^{\nu}}^{2^{\nu-i-1}}) \right). 
\end{gather*}
Here $i\in \{0,\ldots,\nu-1\}$. We denote by $\Phi_{m,\nu,\ast}\colon R_{m,\nu,\ast}\rightarrow I_{m,\nu}$ the sum of $\psi_{m,\nu,i}$ for all $i\in \{0,1,2\}$ and $\omega_{m,\nu,0}$. Then there is an isomorphism $\Coker(\Phi_{m,\nu,\ast})\cong (\Z/m\Z)^{\oplus 2^{\nu-1}-1}$, which follows from the following: 
\begin{gather*}
\left(1,\sum_{j=0}^{(m-3)/2}\rho_{m}^{1+2j}(1+\sigma_{2^{\nu}}^{-1})-\sum_{i=0}^{m-1}\rho_{m}^{i}\right)=\psi_{m,\nu,0}(1)-\sum_{j=0}^{(m-3)/2}\psi_{m,\nu,2}(\rho_{m}^{1+2j});\\
(0,1-\rho_{m})=\omega_{m,\nu,0}((1+\rho_{m}^{2}+\cdots+\rho_{m}^{m-1})(1+\rho_{m}\sigma_{2^{\nu}}^{2^{\nu-1}}));\\
(0,1-\sigma_{2^{\nu}}^{2^{\nu-1}})=\omega_{m,\nu,0}(1+\rho_{m}\sigma_{2^{\nu}}^{2^{\nu-1}}+\cdots +(\rho_{m}\sigma_{2^{\nu}}^{2^{\nu-1}})^{m-1});\\
m(0,1-\sigma_{2^{\nu}}^{-1})=\psi_{m,\nu,0}(\tau_{2^{\nu}}-1)+(1-\sigma_{2^{\nu}}^{-1})\sum_{i=1}^{m-1}(0,1-\rho_{m}^{i}). 
\end{gather*}
In particular, $\Phi_{m,1,\ast}$ is surjective. On the other hand, for $\nu \geq 2$, we write $\Phi_{m,\nu,\bullet}\colon R'_{m,\nu,\bullet}\rightarrow I_{m,\nu}$ for the sum of $\omega_{m,\nu,i}$ for all $i\in \{1,\ldots,\nu-1\}$. Moreover, $\Phi_{m,\nu} \colon R_{m,\nu}\rightarrow I_{m,\nu}$ denotes the sum of $\Phi_{m,\nu,\ast}$ and $\Phi_{m,\nu,\bullet}$. Then, we have
\begin{equation*}
2^{i}(0,1-\rho_{m}\sigma_{2^{\nu}}^{2^{\nu-i-1}})=\omega_{m,\nu,i}(1)+\sum_{j=1}^{2^{i}-1}(1-\rho_{m}\sigma_{2^{\nu}}^{2^{\nu-i-1}})(0,1-\sigma_{2^{\nu}}^{2^{\nu-i}j})
\end{equation*}
for any $i\in \{1,\ldots,\nu-1\}$. Hence, by induction, $\Phi_{m,\nu}$ is surjective for all $\nu \geq 2$. Consequently, we obtain an exact sequence of $G$-lattices
\begin{equation*}
0\rightarrow U_{m,\nu} \rightarrow R_{m,\nu} \rightarrow I_{m,\nu} \rightarrow 0. 
\end{equation*}
In the following, we write $C_{m}$ and $Z_{\nu}$ for the subgroups of $G_{m,\nu}$ generated by $\rho_{m}$ and $\sigma_{2^{\nu}}^{2^{\nu-1}}$ respectively. In addition, we define $G_{m,\nu}$-sublattices of $R_{m,\nu}$ as follows: 
\begin{equation*}
R_{m,\nu,\ast}^{-}:=R_{m,\nu,0}\oplus R_{m,\nu,1}\oplus R_{m,\nu,2},\quad R_{m,\nu}^{-}:=R_{m,\nu,\ast}^{-}\oplus R'_{m,\nu,\bullet}, 
\end{equation*}

\begin{claim}
\begin{enumerate}
\item There exist isomorphisms of $G_{m,\nu}$-lattices
\begin{align}
U_{m,\nu} \cap R_{m,\nu,\ast}^{-}&\cong \Z[G_{m,\nu}/\langle \rho_{m},\sigma_{2^{\nu}}\rangle]\oplus R_{m,\nu,2}^{C_{m}},\label{eq:gmnr}\\
U_{m,\nu} \cap R_{m,\nu,\ast}&\cong (U_{m,\nu} \cap R_{m,\nu,\ast}^{-}) \oplus (R'_{m,\nu,0})^{C_{m}}\label{eq:gmna}. 
\end{align}
Moreover, the following diagram is commutative: 
\begin{equation}\label{eq:fcua}
\begin{gathered}
\xymatrix@C=35pt{
U_{m,\nu} \cap R_{m,\nu,\ast}^{-} 
\ar[r]^{x\mapsto (x,0)\hspace{40pt}}\ar[d]^{x\mapsto x}& 
(U_{m,\nu} \cap R_{m,\nu,\ast}^{-})\oplus (R'_{m,\nu,0})^{C_{m}}\ar[d]&
(R'_{m,\nu,0})^{C_{m}}
\ar[l]_{\hspace{45pt}x\mapsto (0,x)}\ar[d]^{x\mapsto x}\\
R_{m,\nu,\ast}^{-}\ar[r]^{x\mapsto (x,0)}& 
R_{m,\nu,\ast}^{-}\oplus R'_{m,\nu,0}\ar[r]^{\pr_{2}}&
R'_{m,\nu,0}. 
}
\end{gathered}
\end{equation}
Here, the central vertical map is induced by \eqref{eq:gmna}. 
\item We have $U_{m,\nu}=(U_{m,\nu}\cap R_{m,\nu,\ast})+U_{m,\nu}^{Z_{\nu}}$. 
\end{enumerate}
\end{claim}

\begin{proof}[Proof of Claim]
(i): By the definition of $\Phi_{m,\nu,\ast}$, we have a commutative diagram
\begin{equation*}
\xymatrix@C=35pt{
0\ar[r]& U_{m,\nu}\cap R_{m,\nu,\ast}^{C_{m}} \ar[r] \ar[d] & R_{m,\nu,\ast}^{C_{m}} \ar[r]^{\Phi_{m,\nu,\ast}}\ar[d]& I_{m,\nu}^{C_{m}} \ar[r]\ar[d]& 0\\
0\ar[r]& U_{m,\nu}\cap R_{m,\nu,\ast} \ar[r] & R_{m,\nu,\ast} \ar[r]^{\Phi_{m,\nu,\ast}} & I_{m,\nu},&
}
\end{equation*}
where the horizontal sequences are exact. Note that the central and the rightmost vertical maps are injective, and the cokernel of the central vertical map is torsion-free. Moreover, since the cokernel of $\Phi_{m,\nu,\ast}$ is torsion, we have 
\begin{align*}
\rk_{\Z}(U_{m,\nu}\cap R_{m,\nu,\ast})&=\rk_{\Z}(R_{m,\nu,\ast})-\rk_{\Z}(I_{m,\nu})\\
&=(2^{\nu+1}(m+1)+1)-(2^{\nu+1}m-1)\\
&=2^{\nu+1}+2. 
\end{align*}
Therefore, the ranks of $U_{m,\nu}\cap R_{m,\nu,\ast}$ and $U_{m,\nu}\cap R_{m,\nu,\ast}^{C_{m}}$ coincide. Hence, Lemma \ref{lem:lmis} implies
\begin{equation*}
U_{m,\nu}\cap R_{m,\nu,\ast}^{C_{m}}=U_{m,\nu}\cap R_{m,\nu,\ast}. 
\end{equation*}
On the other hand, Theorem \ref{thm:dnqp} gives an exact sequence of $G_{m,\nu}$-lattices
\begin{equation*}
0\rightarrow \Z[G_{m,\nu}/\langle \rho_{m},\sigma_{2^{\nu}} \rangle]\rightarrow R_{m,\nu,0}\oplus R_{m,\nu,1} \rightarrow I_{G_{m,\nu}/\cH_{m,\nu}}^{C_{m}}\rightarrow 0. 
\end{equation*}
Hence the assertion follows from Lemma \ref{lem:cfad}. 

(ii): By (i), $U_{m,\nu}\cap R_{m,\nu,\ast}=\Ker(\Phi_{m,\nu,\ast})$ is a permutation $G_{m,\nu}$-lattice. In particular, we have $H^{1}(Z_{\nu},U_{m,\nu}\cap R_{m,\nu,\ast})=0$. Moreover, one has $(R'_{m,\nu,\bullet})^{Z_{\nu}}=R'_{m,\nu,\bullet}$ by definition. Hence, the assertion follows from Lemma \ref{lem:fgdc} for the commutative diagram
\begin{equation*}
\xymatrix{
0\ar[r]& R_{m,\nu,\ast} \ar[r] \ar[d]^{\Phi_{m,\nu,\ast}} & R_{m,\nu} \ar[r]^{\pi'_{m,\nu}}\ar[d]^{\Phi_{m,\nu}}& R'_{m,\nu,\bullet} \ar[r]\ar[d]& 0\\
0\ar[r]& I_{m,\nu} \ar@{=}[r] & I_{m,\nu} \ar[r] & 0 \ar[r]&0, 
}
\end{equation*}
where $\pi'_{m,\nu}$ is the canonical projection $R_{m,\nu}\twoheadrightarrow R'_{m,\nu,\bullet}$. 
\end{proof}

\vspace{6pt}

In the following, we prove that there is an isomorphism of $G_{m,\nu}$-lattices
\begin{equation}\label{eq:fmns}
U_{m,\nu}\oplus S_{m,\nu} \cong (U_{m,\nu}\cap R_{m,\nu,\ast}^{-})\oplus S'_{m,\nu}\oplus S_{m,\nu},
\end{equation}
where
\begin{gather*}
S_{m,\nu}:=\bigoplus_{i=1}^{\nu-1}\Z[G_{m,\nu}/\langle \rho_{m},\sigma_{2^{\nu}}^{2^{\nu-i}},\sigma_{2^{\nu}}\tau_{2^{\nu}}\rangle],\quad
S'_{m,\nu}:=(R'_{m,\nu,0})^{C_{m}}\oplus \left(\bigoplus_{i=1}^{\nu-1}\Z[G_{m,\nu}/\langle \sigma_{2^{\nu}}^{2^{\nu-i}},\sigma_{2^{\nu}}\tau_{2^{\nu}}\rangle]\right),
\end{gather*}
such that the following diagram is commutative: 
\begin{equation*}
\xymatrix{
U_{m,\nu}\cap R_{m,\nu,\ast}^{-}\ar[r]^{x\mapsto x}\ar[d]^{x\mapsto (x,0,0)}&
U_{m,\nu}\ar[d]^{x\mapsto (x,0)}\\
(U_{m,\nu}\cap R_{m,\nu,\ast}^{-})\oplus S'_{m,\nu}\oplus S_{m,\nu}
\ar[r]^{\hspace{40pt}\cong}&
U_{m,\nu}\oplus S_{m,\nu}. 
}
\end{equation*}
This implies the desired assertion, since \eqref{eq:gmnr} implies that $U_{m,\nu}\cap R_{m,\nu,\ast}^{-}$ is permutation. We prove \eqref{eq:fmns} by induction on $\nu$. If $\nu=1$, then the $G_{m,\nu}$-lattice $R'_{m,\nu,\bullet}$ is zero. Hence, the assertion follows from Claim (i). Next, let $\nu \geq 2$, and suppose that the assertion holds for $\nu-1$. By definition, there is a canonical isomorphism $I_{m,\nu-1}\cong I_{m,\nu}^{Z_{\nu}}$. Moreover, the restriction of $\Phi_{m,\nu}$ to $R_{m,\nu}^{-}$ induces a coflabby resolution of $I_{m,\nu}^{Z_{\nu}}$ as follows: 
\begin{equation*}
0\rightarrow U_{m,\nu}\cap (R_{m,\nu}^{-})^{Z_{\nu}} \rightarrow (R_{m,\nu}^{-})^{Z_{\nu}} \rightarrow I_{m,\nu}^{Z_{\nu}}\rightarrow 0. 
\end{equation*}
Consequently, Lemma \ref{lem:cfad} gives an isomorphism
\begin{equation*}
U_{m,\nu}^{Z_{\nu}} \cong (U_{m,\nu}\cap (R_{m,\nu}^{-})^{Z_{\nu}})\oplus (R'_{m,\nu,0})^{Z_{\nu}}
\end{equation*}
and a commutative diagram
\begin{equation}\label{eq:xitc}
\begin{gathered}
\xymatrix@C=35pt{
U_{m,\nu}\cap (R_{m,\nu}^{-})^{Z_{\nu}}
\ar[r]^{x\mapsto (x,0)\hspace{35pt}}\ar[d]^{x\mapsto x}&
(U_{m,\nu}\cap (R_{m,\nu}^{-})^{Z_{\nu}})\oplus (R'_{m,\nu,0})^{Z_{\nu}}\ar[d]&
(R'_{m,\nu,0})^{Z_{\nu}} \ar[l]_{\hspace{55pt}x\mapsto (0,x)}\ar[d]^{x\mapsto x}\\
R_{m,\nu}^{-}\ar[r]^{x\mapsto (x,0)}& R_{m,\nu}^{-}\oplus R'_{m,\nu,0} \ar[r]^{\pr_{2}}&
R'_{m,\nu,0}. 
}
\end{gathered}
\end{equation}
On the other hand, by the induction hypothesis, we obtain an isomorphism of $G_{m,\nu}$-lattices
\begin{equation}\label{eq:fsdd}
(U_{m,\nu}\cap (R_{m,\nu}^{-})^{Z_{\nu}}) \oplus S_{m,\nu-1} \cong (U_{m,\nu}\cap (R_{m,\nu,\ast}^{-})^{Z_{\nu}})\oplus S'_{m,\nu-1} \oplus S_{m,\nu-1}
\end{equation}
such that the following diagram is commutative: 
\begin{equation*}
\xymatrix{
U_{m,\nu}\cap (R_{m,\nu,\ast}^{-})^{Z_{\nu}}\ar[r]^{x\mapsto x}\ar[d]^{x\mapsto (x,0,0)}&
U_{m,\nu}\cap (R_{m,\nu}^{-})^{Z_{\nu}}\ar[d]^{x\mapsto (x,0)}\\
(U_{m,\nu}\cap (R_{m,\nu,\ast}^{-})^{Z_{\nu}})\oplus S'_{m,\nu-1} \oplus S_{m,\nu-1}\ar[r]^{\hspace{20pt}\cong}&
(U_{m,\nu}\cap (R_{m,\nu}^{-})^{Z_{\nu}}) \oplus S_{m,\nu-1}. 
}
\end{equation*}

Now, we denote by $\Xi_{m,\nu,1}$ the central vertical homomorphism in \eqref{eq:fcua}, that is, 
\begin{equation*}
    (U_{m,\nu}\cap R_{m,\nu,\ast}^{-})\oplus (R'_{m,\nu,0})^{C_{m}} \rightarrow R_{m,\nu,\ast}^{-}\oplus R'_{m,\nu,0}
\end{equation*}
On the other hand, set $S_{m,\nu-1}^{\dagger}:=S'_{m,\nu-1}\oplus S_{m,\nu-1}$, and consider the $G$-homomorphism
\begin{equation*}
\Xi_{m,\nu,2}\colon S_{m,\nu-1}^{\dagger}\oplus (R'_{m,\nu,0})^{Z_{\nu}}\rightarrow R_{m,\nu}\oplus S_{m,\nu-1}=R_{m,\nu,\ast}^{-}\oplus R'_{m,\nu,0}\oplus (R'_{m,\nu,\bullet}\oplus S_{m,\nu-1}),
\end{equation*}
induced by \eqref{eq:fsdd} and the central vertical map of \eqref{eq:xitc}. Denote by $\Xi_{m,\nu}$ the sum of $\Xi_{m,\nu,1}$ and $\Xi_{m,\nu,2}$. Then we have
\begin{equation}\label{eq:imxi}
    \Ima(\Xi_{m,\nu})=U_{m,\nu}\oplus S_{m,\nu-1},
\end{equation}
which is a consequence of Claim (ii) and \eqref{eq:fsdd}. In particular, one has an exact sequence
\begin{equation}\label{eq:ffex}
0\rightarrow E_{m,\nu}\rightarrow (U_{m,\nu}\cap R_{m,\nu,\ast}) \oplus S_{m,\nu-1}^{\dagger} \oplus (R'_{m,\nu,0})^{Z_{\nu}} \xrightarrow{\Xi_{m,\nu}} U_{m,\nu}\oplus S_{m,\nu-1}\rightarrow 0. 
\end{equation}

Now, we confirm that $\Xi_{m,\nu,1}$ and $\Xi_{m,\nu,2}$ satisfy the four assumptions in Lemma \ref{lem:cssp}. The condition (1) is not difficult by using the property that $m$ is odd. Moreover, (2) is a consequence of Claim (i) and \eqref{eq:xitc}. On the other hand, \eqref{eq:fsdd} induces the following: 
\begin{align*}
\Xi_{m,\nu,1}(U_{m,\nu}\cap (R_{m,\nu,\ast}^{-})^{Z_{\nu}})+\Xi_{m,\nu,2}(S_{m,\nu-1}^{\dagger})&=(U_{m,\nu}\cap (R_{m,\nu}^{-})^{Z_{\nu}})\oplus S_{m,\nu-1}\\
&=(U_{m,\nu}\oplus S_{m,\nu-1})\cap (R_{m,\nu}^{-}\oplus S_{m,\nu-1})^{Z_{\nu}}. 
\end{align*}
Hence, (3) is valid by \eqref{eq:imxi}. Finally, \eqref{eq:ffex} implies $\rk_{\Z}(\Ker(\Xi_{m,\nu}))=2^{\nu-1}$, that is, (4) holds. Therefore, we can apply Lemma \ref{lem:cssp}. In particular, we obtain an isomorphism of $G_{m,\nu}$-lattices
\begin{equation}\label{eq:aflm}
\begin{gathered}
U_{m,\nu}\oplus S_{m,\nu-1}\oplus \Z[G_{m,\nu}/H_{\nu}C_{m}Z_{\nu}]
\\ \cong (U_{m,\nu}\cap R_{m,\nu,\ast}^{-})\oplus S_{m,\nu-1}^{\dagger}\oplus \Z[G_{m,\nu}/H_{\nu}C_{m}]\oplus \Z[G_{m,\nu}/H_{\nu}Z_{\nu}]
\end{gathered}
\end{equation}
where $H_{\nu}:=\langle \sigma_{2^{\nu}}\tau_{2^{\nu}}\rangle$, which induces the restriction of $\Xi_{m,\nu}$ to $(U_{m,\nu}\cap R_{m,\nu,\ast}^{-})\oplus S_{m,\nu-1}^{\dagger}$. On the other hand, by the definitions of $S_{m,\nu}$ and $S'_{m,\nu}$, there exist isomorphisms
\begin{gather*}
S_{m,\nu-1}\oplus \Z[G_{m,\nu}/H_{\nu}C_{m}Z_{\nu}]=S_{m,\nu},\\
S_{m,\nu-1}^{\dagger}\oplus \Z[G_{m,\nu}/H_{\nu}C_{m}]\oplus \Z[G_{m,\nu}/H_{\nu}Z_{\nu}]
\cong S'_{m,\nu}\oplus S_{m,\nu}.
\end{gather*}
Consequently, \eqref{eq:aflm} can be written as \eqref{eq:fmns}. Hence, the proof is complete. 
\end{proof}

\begin{thebibliography}{99}
\bibitem[End11]{Endo2011} S.~Endo, \emph{The rationality problem for norm one tori}, Nagoya Math.~J.~\textbf{202} (2011), 83--106.
\bibitem[Lor05]{Lorenz2005} M.~Lorenz, \emph{Multiplicative Invariant Theory}, Encyclopaedia of Mathematical Sciences, vol.~135, Springer, 2005.
\end{thebibliography}
\end{document}
